\documentclass[instructions.tex]{subfiles}
\begin{document}
 
\chapter{Il faut, pendant la vie, apprendre à bien mourir.}
 
\primo Vous appréhendez la mort; vous avez raison de l'appréhender si vous ne vivez pas en chrétien. Comment ne pas craindre la mort quand on ne s'y prépare pas, et que l'on n'a jamais appris à mourir? Il est bien tard au dernier moment d'apprendre à faire une chose qui demande une préparation de toute la vie.
 
Les saints, sur la terre, ne faisaient autre chose que d'apprendre à bien mourir, et craignaient toujours de ne pas réussir. J'apprends chaque jour à mourir, disait saint Paul\footnote{Quotidiè morior. 0r. 15. 31.}. Le cardinal Bellarmin, après avoir vécu dans la sainteté, demandait à Dieu, à l'heure de la mort, encore une année pour se disposer à ce terrible moment.
 
Un saint homme, ayant passé dans le désert plus de quarante ans dans la pénitence, disait qu'il n'avait fait, pendant tout ce temps, que d'apprendre à mourir. C'est là en effet l'unique chose qu'on devrait avoir à cœur, puisque la vie ne nous est donnée que pour apprendre à bien la finir. A quoi sert l'étude des sciences profanes, et les mouvemens qu'on se donne pour se faire un grand nom, pour augmenter sa fortune et sa réputation, si l'on s'oublie soi-même? Notre vie est un voyage pour l'éternité; elle passe comme une ombre: c'est donc être bien insensé, dit saint Eucher, de s'amuser à tant de choses inutiles pendant un temps si court, tandis qu'on oublie le terme qui doit décider de tout.
 
Noé travailla cent ans à bâtir l'arche pour se préserver du déluge d'eau. Vous n'avez pas cent ans à vivre, vous qui, peut-être avant trois jours, serez dans le tombeau; et cependant, loin de penser à votre dernière heure pour vous préserver d'un déluge de feu, à peine y pensez-vous une fois dans l'année. Un déluge de feu, un enfer, qui est la suite d'une mauvaise mort, est-il donc une chose qui mérite si peu vos réflexions?
 
\secundo Apprenez tous les jours à mourir; pensez chaque jour que la mort est à vos côtés; qu'elle vous arrivera bientôt, peut-être demain; pensez-y sérieusement. On ne meurt qu'une fois: une mauvaise mort est un mal sans remède. Ce n'est pas assez d'y penser, il faut être toujours sur vos gardes, il faut même, dit Jésus-Christ, à chaque moment être prêt, parce que le Seigneur viendra quand vous n'y penserez pas. Ceux qui vous ont précédé ont été surpris, vous le serez vous-même; leur mort inopinée est un présage de la vôtre.
 
Si vous étiez menacé d'être servi, dans un repas, d'une viande empoisonnée, vous seriez toujours en garde contre ce qu'on vous présenterait. A tous les momens de votre vie vous êtes menacé de la mort; il y a un de ces momens qui sera le dernier, qui est connu de Dieu seul; pourquoi ne les appréhendez-vous pas tous? pourquoi n'êtes-vous pas sur vos gardes, en vivant comme si à chaque moment Dieu devait vous appeler? Celui-là est indigne de vivre et de porter le nom de chrétien, dit saint Jérôme, qui a l'assurance de vivre dans un état dans lequel il ne voudrait pas mourir\footnote{Indignus est nomine christiani, qui in eo statu vult vivere, in quo nollet mori. S. Hier.}.
 
Oh! si tous les jours, comme saint Antoine, on avait à son réveil cette pensée : Ce jour est peut-être le dernier de ma vie! si chacun se regardait à tous momens comme un criminel à qui l'on va prononcer la sentence, serait-il possible qu'on commît le péché? y aurait-il un jour qu'on ne vécût comme on voudrait mourir? méditez cette pensée tous les jours, puisque chaque jour vous pouvez mourir.
 
\section{Résolutions}
 
\primo Tous les matins, en me levant, je me dirai à moi-même : Ce jour est peut-être le dernier de ma vie; il faut donc que je le passe en me disposant sérieusement à la mort.
 
\secundo Tous les soirs, je ferai un acte de contrition pour me préparer à ce moment si critique. 

\prayer{O mon Dieu! aidez-moi à mettre ordre au passé, et à profiter du présent, de telle sorte que la mort me trouve prêt à quitter saintement la vie.}
 
\end{document}
